\documentclass[11pt]{article}
 
% Latin Modern: the modern successor to Computer Modern
\usepackage[T1]{fontenc}
\usepackage{lmodern}
\usepackage[utf8]{inputenc}
 
\usepackage[margin=1in]{geometry}
\usepackage{parskip}      % space between paragraphs, no indent
\usepackage{enumitem}
\usepackage{setspace}
\singlespacing            % single-spaced body
 
\pagestyle{plain}
 
\begin{document}
 
% --- Header (top left) ---
\noindent
Daniel Kadyrov\\
\today\\
200-Hour Teacher Training 2026-2027\\
Sonic Yoga Center
 
\vspace{1em}
 
% --- Title: centered, body-text size, bold ---
\begin{center}
  \textbf{Karma Yoga Project}
\end{center}
 
One of the Sonic Yoga Teaching Training Teachers, Mitch Hipp, often ends his classes with the phrase: ``What if everyone in the world practiced yoga? My bet is that the world would be a better place.'' After a physically and mentally engaging asana class, this phrase often lives in my mind well past my mat, as a bicycle home, make my meals, and before I take my deep yoga nap (unless it's his \textit{jetlag class} in which it's before I pass out for my nightly sleep). 

My current employment is as an engineer in a small engineering federal laboratory, a small, and what sometimes seems like an insignificant part, of the currently problematic Department of Homeland Security. Although I find my coworkers to be in similar social, moral, and political views as myself and the general New York City population, my office is situated in a federal building that is also home to agencies including Immigration and Customs Enforcement (ICE), Customs and Border Protection (CBP), and Homeland Security Investigation (HSI), groups that have recently become notable in their aggressive enforcement of the Trump Administration policies. 

Although, by recommendation of the Sonic Yoga Teacher Training Teachers, I already started teaching to my coworkers in preparation for the final practicum, my Karma Yoga project once I have my certification is to formally continue teaching, as a free offering, to my fellow federal employees that work at the 201 Varick Federal Building. Beyond bringing physical movement, strength and mobility, to workers who often sit for hours at their desk jobs, this offering would extend to federal employees including those in the more physically demanding roles which are often in a demographic that have unfortunately misaligned views towards yoga. 

In this way, I hope to slowly try to fulfill my teacher's wish by bringing yoga to a population that might need their hamstrings stretched, get humbled in a hand balance, feel the emotional release of a hip opener, get the body awareness of their chakras, and maybe put their hands in Anjali mudra, chant Ohm, and close their practice with Lokah Samastah Sukhino Bhavantu.

% Create an Intention for a Karma Yoga Project to be implemented post graduation. Type a one-paragraph explanation of your intention to be handed in.

% The Karma Yoga Project is a typed paragraph explaining how you want to serve the community with yoga. At the end of the day, we are giving Yoga. You can write about volunteering for a program that you care deeply about, or offering Yoga to people who need to heal. It does not have to specifically be teaching a free Yoga asana class, but the knowledge you have accumulated from your training and Yoga practice should in some way inform the work you do. Find something that sparks your interest and ignites the light inside!

% Questions to consider when writing your karma yoga paragraph:
% What do you want to do with your Yoga?
% How do you want to give back to the community?
 
\end{document}